\documentclass[12pt,a4paper]{article}
\usepackage[utf8]{inputenc}
\usepackage[T2A]{fontenc}
\usepackage[russian]{babel}
\usepackage{amsmath,amssymb}
\usepackage{geometry}
\usepackage{hyperref}
\geometry{margin=2.5cm}

\title{Проверка трёх вариантов механизма электрического кулона \\
для модели прецессирующих магнитных диполей Дроздова}
\author{Расчёты выполнены совместно с (Claude Sonnet 5)}
\date{\today}

\begin{document}
\maketitle

\begin{abstract}
В работе проверяются три варианта механизма возникновения эффективного
электрического кулоновского взаимодействия между двумя нейтральными
(по суммарному магнитному заряду) прецессирующими магнитными диполями
Дирака -- моделью электрона, предложенной А.~Дроздовым. Показано, что
линейная теория (Вариант 1, классическая магнитостатика Пича--Кёлера) и
линейная теория с корректной геометрией границы Ферми (Вариант 2) дают
только быстро затухающие ($\propto R^{-3}$) или безнаковые ($\propto R^{-1}$,
но не зависящие от разности фаз прецессии) вклады. Вариант 3, основанный
на гипотезе о нелинейной (Янг--Миллс-подобной) самосвязи калибровочного
поля дислокаций, впервые даёт знаковый, не гасящийся при усреднении по
прецессии вклад $\propto R^{-1}\cos\Delta\phi$ -- то есть верную
\emph{функциональную форму} кулоновского закона. Однако величина эффекта
при дименсионно естественных оценках константы самосвязи $\lambda$
отличается от наблюдаемого кулона на 5--10 порядков.
\end{abstract}

\tableofcontents

\section{Постановка задачи}

Рассматривается система из двух прецессирующих магнитных диполей
(``электронов'' в модели Дроздова), каждый состоящий из пары монополей
Дирака $+q_m,-q_m$, вращающихся с частотой $\omega$ по окружности радиуса
$a$. Диполь 1 центрирован в начале координат, диполь 2 -- на расстоянии
$R$ вдоль оси $x$. Разность фаз прецессии обозначена $\Delta\phi=\phi_1-\phi_2$.

Цель -- проверить, может ли взаимодействие таких (в сумме электрически
нейтральных по магнитному заряду) объектов дать эффективный электрический
кулон $U\propto 1/R$ со знаком, зависящим от взаимной ориентации диполей.

\section{Вариант 1: линейная магнитостатика (старая физика)}

Основа -- теория Кадича--Эделена, на которую опирается расчёт
Брагинского. Эта теория явно линейна: калибровочная теория
дислокаций Кадича--Эделена не содержит калибровочного взаимодействия,
в исходных работах была построена линейная теория, без взаимодействия.

Знаковая сумма потенциалов по всем четырём парам монополей:
\begin{equation}
U_{signed}(R,\phi_1,\phi_2) = q_m^2\left(
\frac{1}{r_{++}} - \frac{1}{r_{+-}} - \frac{1}{r_{-+}} + \frac{1}{r_{--}}
\right)
\end{equation}
Разложение по малому параметру $\varepsilon=a/R$ дает точный результат:
\begin{equation}
U_{signed}\cdot R = q_m^2\Big[-2\varepsilon^2\cos(\phi_1-\phi_2)
-6\varepsilon^2\cos(\phi_1+\phi_2)\Big] + O(\varepsilon^4)
\end{equation}
Члены нулевого ($R^{-1}$, монополь) и первого ($R^{-2}$, диполь) порядка
отсутствуют тождественно -- это прямое следствие теоремы Гаусса:
суммарный магнитный заряд каждого диполя $q_m-q_m=0$ точно в любой момент
времени. Усреднение по прецессии ($\phi_1\to\phi_1+\omega t$) дает
единственный не исчезающий член:
\begin{equation}
\langle U_{signed}\rangle \cdot R = -2q_m^2\varepsilon^2\cos\Delta\phi
+ O(\varepsilon^4)
\end{equation}

\textbf{Вывод.} Ведущий вклад -- квадрупольный, $\propto R^{-3}$.
Настоящего кулона $R^{-1}$ данный (линейный) механизм не дает.

\section{Вариант 2: линейная плёнка + геометрия Ферми (Вариация B)}

При корректной релятивистски инвариантной границе (условие Ферми:
нормаль к плёнке пропорциональна 4-скорости трубки) и линейной
минимизации площади плёнки, источник продавливания вводится как
\begin{equation}
\sigma_i = \frac{m_i\,(\vec W_i\cdot\vec W_i)}{c^2},
\end{equation}
где $\vec W_i$ -- собственное (центростремительное) ускорение
прецессии. Так как $\sigma_i\ge 0$ независимо от знака заряда,
результирующее поле -- истинный монополь:
\begin{equation}
\delta w(R) = -\frac{\sigma_i}{\gamma_{3D}\,R},\qquad
\gamma_{3D}=\frac{c^7}{8\pi\hbar G^2}
\end{equation}
Этот вклад дает $R^{-1}$-хвост, но он не зависит от знака
заряда и от разности фаз $\Delta\phi$ -- то есть он гравитационно-,
а не электрически-подобный.

Численная проверка (параметры электрона, $v=0.03c$--$0.09c$
из условия магнетона Бора): расхождение с реальной величиной
$Gm_e^2/e^2\approx 2.4\times10^{-43}$ составляет $10^{40}$--$10^{50}$
раз в зависимости от способа калибровки жесткости $\gamma_{3D}$.

\textbf{Вывод.} Ни знака, ни нужной величины эффект не воспроизводит.

\section{Вариант 3: нелинейная (Янг--Миллс-подобная) самосвязь плёнки}

\subsection{Физическое обоснование нелинейности}

Калибровочная теория дислокаций в её полной форме нелинейна: полевые
уравнения имеют структуру типа Янга--Миллса, и из-за сложности
калибровочного взаимодействия возникают самозацепления, связывающие
дислокационный импульс и дислокационное напряжение. Точный вид членов
самосвязи в доступных материалах не восстановлен; ниже используется
модельная (но физически мотивированная) кубическая нелинейность
\begin{equation}
\nabla^2\Phi = -4\pi\rho_{lin} - \lambda\,\Phi^2,
\end{equation}
где $\rho_{lin}$ -- линейная (дипольная) плотность заряда, а $\lambda$ --
константа самосвязи размерности [см$^{-2}$].

\subsection{Аналитический расчёт крест-члена}

Возмущение первого порядка по $\lambda$ в дальней зоне определяется
интегралом
\begin{equation}
\Phi_1(R\to\infty) \approx \frac{\lambda}{4\pi R}\int \Phi_0^2\,dV.
\end{equation}
Для взаимодействия двух диполей интересен именно крест-член
$I(R)=\int \Phi_{0,1}\Phi_{0,2}\,dV$, а не самоэнергия $\int\Phi_0^2\,dV$
(которая $R$-независима). Через фурье-свёртку двух точечных диполей
с моментами $\vec p_1,\vec p_2$ и разделением $\vec R$ получена
закрытая форма (проверена численно методом Монте-Карло с точностью
около 5\%):
\begin{equation}
I(R) = \frac{2\pi}{R}\Big[\vec p_1\cdot\vec p_2
- (\vec p_1\cdot\hat R)(\vec p_2\cdot\hat R)\Big].
\end{equation}
Для прецессирующих диполей ($\vec p_i = p(\cos\phi_i,\sin\phi_i,0)$,
$\vec R\parallel \hat x$) и усреднения по времени:
\begin{equation}
\boxed{\ \langle I(R,\Delta\phi)\rangle
= \frac{\pi p^2}{R}\cos\Delta\phi,\qquad p=2q_m a\ }
\end{equation}

Это -- знаковый, не гасящийся при усреднении по прецессии,
$\propto R^{-1}$ эффект. Структурно это первый из проверенных
механизмов, дающий верную функциональную форму электрического кулона.

\subsection{Оценка величины эффекта}

Энергия нелинейного взаимодействия оценивается как
\begin{equation}
U_{nl}(R) \approx \lambda\,\frac{\pi p^2}{R}\cos\Delta\phi.
\end{equation}
При дименсионно естественной оценке $\lambda\sim 1/r_0^2$
($r_0$ -- классический радиус электрона):
\begin{equation}
\frac{U_{nl}(a_0)}{U_{Coulomb}(a_0)} \approx 4.5\times10^{9}
\end{equation}
При более консервативной оценке $\lambda\sim 1/a^2$
($a$ -- комптоновская длина волны):
\begin{equation}
\frac{U_{nl}(a_0)}{U_{Coulomb}(a_0)} \approx 6.0\times10^{4}
\end{equation}
Точное совпадение с реальным кулоном требует
\begin{equation}
\lambda_{needed} \approx 1.1\times10^{16}\ \text{см}^{-2}
\approx 2.2\times10^{-10}\cdot\frac{1}{r_0^2}.
\end{equation}

\section{Сводная таблица результатов}

\begin{center}
\begin{tabular}{|l|c|c|c|}
\hline
Вариант & Ведущий порядок & Знак/$\Delta\phi$? & Кулон $1/R$ со знаком? \\
\hline
1. Линейная магнитостатика & $R^{-3}$ & да ($\cos\Delta\phi$) & Нет \\
2. Линейная плёнка Ферми & $R^{-1}$ & нет & Нет (нет знака) \\
3. Нелинейная самосвязь & $R^{-1}$ & да ($\cos\Delta\phi$) & Да, но $\lambda$ не откалибрована \\
\hline
\end{tabular}
\end{center}

\section{Заключение}

Вариант 3 -- единственный из проверенных механизмов, дающий правильную
функциональную форму электрического кулона: знаковую, линейную по
$1/R$, зависящую от взаимной фазы прецессии. Величина эффекта, однако,
целиком определяется константой самосвязи $\lambda$, точное значение
которой не выведено из первых принципов калибровочной теории дислокаций
в доступных материалах, а лишь оценено по аналогии со структурой
самозацепления Янга--Миллса. При естественных оценках $\lambda$
эффект получается на 5--10 порядков больше наблюдаемого кулона -- то есть
свобода произвола, ранее находившаяся в показателе степени $R$ или в
показателе степени массы, переместилась в величину $\lambda$, но не
исчезла.

\begin{thebibliography}{9}
\bibitem{Fermi1923} E. Fermi, {\it О парадоксе массы в электродинамике},
Phys. Z. {\bf 24}, 340 (1923).
\bibitem{Braginsky2018} А. Я. Брагинский и др., материалы Braginsky\_12\_2018.pdf.
\bibitem{Lazar2010} M. Lazar, {\it Gauge theory of dislocations},
lazar2010.pdf (полный текст не восстановлен из имеющихся материалов --
доступны лишь метаданные PDF-профиля печати).
\bibitem{Braginsky2016} А. Я. Брагинский, Braginsky\_2016\_.doc.
\end{thebibliography}

\end{document}
